\subsection{The exact tame configuration of a possible quintic class obstruction}

\begin{proposition}[KT-QUINTIC-TAME-CONFIGURATION-LIMIT]
Let $p\equiv1\pmod {10}$ be prime, let $F$ be the cyclic quintic
subfield of $\mathbf Q(\mu_p)$, and assume the remaining possibility
allowed by KT-ORDER-FIVE-LENGTH-ONE:
$A=\operatorname{Cl}(F)\{p\}\simeq C_p$.
Let $H/F$ be its ordinary Hilbert $p$-class field.  Then
\[
 \operatorname{Gal}(H/\mathbf Q)=P\rtimes I,
 \qquad P\simeq C_p,\quad I\simeq C_5,                 \tag{QTC.1}
\]
where the action is faithful.  At every prime over $p$, the
decomposition group and inertia group coincide and are conjugate to
$I$.  Consequently the prime of $F$ over $p$ splits completely in
$H/F$.  If $K=H^I$, then
\[
 \begin{gathered}
 [K:\mathbf Q]=p,\qquad
 p\mathcal O_K=\mathfrak q_0
       \prod_{i=1}^{(p-1)/5}\mathfrak q_i^5,
 \qquad f(\mathfrak q_i/p)=1,\\
 d_H=p^{4p},\qquad d_K=p^{4(p-1)/5},\qquad
 \operatorname{rd}(H)=p^{4/5},\quad
 \operatorname{rd}(K)=p^{4(p-1)/(5p)}.                 \tag{QTC.2}
 \end{gathered}
\]
The fields $H$ and $K$ are totally real.

Conversely, any tamely ramified $C_p\rtimes C_5$ Galois extension
of $\mathbf Q$, with faithful action and ramified at no finite prime
other than this same $p$, gives an unramified $C_p$ extension of $F$.
Its decomposition group at $p$ must already be the inertia complement.
Thus the complete-splitting condition contributed by the principal
$p$-prime is not an additional restriction on this tame configuration.
These assertions do not exclude any of the four possible primitive
order-five class characters.
\end{proposition}

\begin{proof}
The Hilbert $p$-class field is preserved by every automorphism of
$F/\mathbf Q$, by its defining maximality and uniqueness.  Its Galois
group therefore fits into $1\to A\to G\to C_5\to1$.
The $C_5$-action on $A$ is nontrivial: if $x\in A$ is fixed, the
norm-extension identity gives
$5x=jN_{F/\mathbf Q}x=0$, because $\operatorname{Cl}(\mathbf Q)=0$.
Since $p\ne5$, this forces $x=0$.  A nontrivial action of $C_5$ on
$C_p$ is faithful.  Cauchy's theorem supplies a subgroup of order
five in $G$; its intersection with $A$ is trivial, so it is a
complement.  Identify the resulting affine group with
\[
 G=\langle x,s:x^p=s^5=1,\ sxs^{-1}=x^\lambda\rangle,
 \qquad \lambda\in\mathbf F_p^\times,
 \quad\operatorname{ord}(\lambda)=5.                 \tag{QTC.3}
\]

The unique prime $\mathfrak p$ of $F$ over $p$ is principal:
the norm to $F$ of $1-\zeta_p$ generates it.  The ordinary class-field
splitting law already implies complete splitting in $H/F$.
There is also a group-theoretic proof that shows why this supplies
no extra constraint.  Since $H/F$ is unramified and $F/\mathbf Q$
is totally ramified of degree five at $p$, inertia in $H/\mathbf Q$
has order five and is a complement.  Every complement has the form
of the stabilizer of one point in the affine action on $\mathbf F_p$.
It is self-normalizing: for $I=\langle s\rangle$, conjugation by
$x^a$ sends its fixed point $0$ to $a$, and $I$ has the unique fixed
point $0$ because $\lambda\ne1$.  Thus $N_G(I)=I$.
Inertia is normal in the decomposition group, whence
$I\subseteq D\subseteq N_G(I)=I$.  This proves $D=I$.

There are $p$ primes of $H$ over $p$, each with completion
$M=F_{\mathfrak p}$, a totally ramified cyclic degree-five
extension of $\mathbf Q_p$.  The coset action on $G/I$ is the
affine action on $\mathbf F_p$.  Its restriction to $I$ has one
fixed point and $(p-1)/5$ free orbits of size five.  The corresponding
completions of $K$ are therefore
\[
 K\otimes_{\mathbf Q}\mathbf Q_p
       \simeq\mathbf Q_p\times M^{(p-1)/5}.           \tag{QTC.4}
\]
For example, this follows directly by taking $I$-invariants in
$H\otimes\mathbf Q_p$, whose factors are indexed by the primes
of $H$ over $p$: the stabilizer of the fixed factor is all of $I$,
and every other factor orbit has trivial stabilizer.
This proves the asserted factorization.

The already used conductor--discriminant formula for $F$ is
$d_F=p^4$.  Since $p$ is its only ramified prime and has a single
completion, the local discriminant of $M/\mathbf Q_p$ is $p^4$.
The trace form of a product of finite extensions is block diagonal,
so local discriminants multiply.  Equation (QTC.4) gives $d_K$,
and the product of $p$ copies of $M$ gives $d_H$.  No other finite
prime is ramified.  Both fields are totally real because the ordinary
Hilbert class field preserves real places; equivalently $G$ has odd
order and hence cannot contain a nontrivial complex conjugation.
Taking the indicated roots proves (QTC.2).

For the converse let $H'/\mathbf Q$ be the specified tame extension.
Its inertia at $p$ has order prime to $p$ and divides $5p$, so has
order one or five.  Order one would make $H'/\mathbf Q$ unramified
at every finite prime, contrary to the discriminant theorem for
$\mathbf Q$.  Thus it has order five, is a complement, and the
same normalizer calculation gives $D=I$.
The fixed field $F'=H'^P$ is cyclic of degree five, ramified only
at $p$ and tamely there.  Its abelian conductor is $p$:
locally the pro-$p$ group $1+p\mathbf Z_p$ maps trivially to its
order-five Galois group, and the character is unramified elsewhere.
There is no infinite conductor because the extension has odd degree.
Global class field theory for $\mathbf Q$, or its ray class description,
therefore identifies $F'$ with the unique quintic subfield $F$ of
$\mathbf Q(\mu_p)$.  Since $I\cap P=1$, the extension $H'/F$
is unramified at $p$, and it is unramified everywhere else, including
the real places.  It is cyclic of degree $p$, so the ordinary class
group of $F$ has a quotient of order $p$.
\end{proof}

\paragraph{The four local labels are all compatible.}
Here is an explicit local check, not a global field construction.
Fix $M=F_{\mathfrak p}$ with its given $C_5$-action.  For each of the
four choices of primitive $\lambda$ in (QTC.3), the induced algebra
\[
 B_\lambda=
 \{f:G_\lambda\to M:f(ig)=i(f(g))\text{ for }i\in I\},
 \qquad (h\cdot f)(g)=f(gh),                          \tag{QTC.5}
\]
is a $G_\lambda$-Galois finite etale $\mathbf Q_p$-algebra.
It is the product of $p$ copies of $M$, and its $P$-invariant
subalgebra is $M$ with the prescribed quotient $C_5$-action.
Evaluation at the identity has stabilizer $I$, which acts on that
factor through the actual Galois action on $M$.
All assertions follow by choosing representatives for $I\backslash G$
in (QTC.5).  In particular, its fixed subalgebra under $G$ is
$M^I=\mathbf Q_p$, and the group acts transitively on the $p$ factors
with the correct factor stabilizer.  This gives the stated Galois
algebra condition.

Thus local tame ramification does not identify the action on the
global normal subgroup $P$ with a distinguished cyclotomic character.
Indeed $P\cap D=1$: these $p$-elements permute completions, rather
than belonging to the local wild inertia group.  Substituting a
local wild-inertia conjugation rule for their action would be an
invalid inference.  The local algebras (QTC.5), and the discriminants
in (QTC.2), are compatible with every primitive label.  They do not
assert any global realization.

\paragraph{What a discriminant exclusion would still need.}
For a totally real degree-$n$ field, Minkowski gives only
\[
 \operatorname{rd}\ge
       \left(\frac{n^n}{n!}\right)^{2/n}<\exp(2).      \tag{QTC.6}
\]
The strict upper bound on the displayed lower-bound expression follows
from $\log(n!)>\int_1^n\log t\,dt=n\log n-n+1$ for $n>1$.
For every $p\ge31$, both root discriminants in (QTC.2) exceed $9$:
for $K$, use $4(p-1)/(5p)>2/3$ and $31^2>9^3$; for $H$
the exponent is still larger.  Thus this bound is compatible with
the exact proposed discriminants.  This comparison is not claimed
as an audit of every refined discriminant bound.

A contradiction by this method needs an additional arithmetic
inequality for the specified tame Frobenius extensions that is
strictly stronger than the exact value in (QTC.2), for example a
lower bound $\operatorname{rd}(K)>p^{4(p-1)/(5p)}$ for that family,
or an independent nonexistence theorem for its fixed ramification
prime $p$.  A realization theorem allowing a different ramified prime
does not supply this input.  Neither the length-one class bound nor
the principal $p$-prime supplies such a strict inequality here.
This bounded attempt removes no class character and is classified
as a limitation of the proposed method, not arithmetic movement.

\paragraph{Exact primary locators and reading scope.}
D.~S.~Dummit and H.~Kisilevsky,
\emph{Decomposition types in minimally tamely ramified extensions
of $\mathbf Q$}, Research in Number Theory \textbf{3}, article~24
(2017), \texttt{doi:10.1007/s40993-017-0085-7}:
the introduction and Definitions~2.1--2.4 on pp.~1--2 give the
inertia/decomposition conventions; Remark~4.1 on p.~18 states
exactly the unique configuration $T=Z=I$ for Frobenius groups with
cyclic complement.  Those passages, the whole p.~18, and the
following splitting explanation were read.  The specialized
normalizer proof is supplied above; no realization theorem at the
specific prime $p$ is imported.

J.~S.~Milne, \emph{Class Field Theory}, version~4.03
(6 August 2020), Theorem~0.4 and the following Hilbert-class-field
paragraph on printed p.~4 give the ordinary class-field splitting
law.  Chapter~V, Theorems~3.5--3.6 and Remark~3.8 on p.~158,
and Examples~3.9--3.10 on p.~159 give the conductor and rational
ray-class statements.  Printed pp.~156--159 and the whole p.~4
were read; the foundational reciprocity and existence proofs are
imported, not freshly reproved.

Milne, \emph{Algebraic Number Theory}, version~3.08
(19 July 2020), Theorem~4.3 on printed p.~70 and Theorem~4.9
with its complete deduction on p.~72 give (QTC.6) and the absence
of a nontrivial everywhere-unramified extension of $\mathbf Q$.
Those pages were read.  The earlier lattice proof of Theorem~4.3
is imported.  The class-character length bound and $d_F=p^4$
are the already checked inputs of KT-ORDER-FIVE-LENGTH-ONE.
The Greenberg, class-pairing, and abstract norm-system notes were
read before this attempt and are not reused as new exclusions.
